\fancyhead[L]{\small SCHULZE ON THE SYSTEM OF THE WORLD.}
\SourceLeaf{538}
\begin{center}\small Göttingische gelehrte Anzeigen. No. 21. pp. 305..306. February 23, 1811.\\[3pt]\rule{.72\linewidth}{.3pt}\end{center}

\emph{An Account of the System of the World. A Guide for Teaching Astronomy in Schools. Written and supplied with the essential bibliographical remarks and references to make it easier to continue studying astronomy on one's own, by} G. L. \textsc{Schulze}, minister at Polenz and Ammelshayn near Leipzig. With 4 copperplates. 1811. 390 pages in octavo. Leipzig, Baumgärtner.

Although the deeper knowledge of astronomy can by its very nature belong only to a few initiates, the science contains such a wealth of highly interesting truths accessible to every educated person that anyone can find in it a lifelong source of varied pleasures and uplifting perspectives. Anyone who had an early opportunity to be shown the way to these truths must always acknowledge it with gratitude. Well-designed instruction in astronomy for young people is therefore very valuable in this respect. There is no doubt that making such instruction more widely available than has been usual among us so far could even awaken many a talent and encourage its further development. In our opinion, the present textbook is very well organized for this purpose. It steers a fortunate middle course between the two reefs on which one otherwise so easily runs aground, the forbidding dryness of excessive abstraction and a shallowness that only fosters half-knowledge and keeps independent thought from taking flight. The references to the finest astronomical writings on every subject treated deserve special praise. Anyone seeking to learn more can satisfy their curiosity there. The purpose of our journal does not allow us to go into more detail about a work of this kind. But we must confess that it was a pleasure for us to find such thorough knowledge of astronomy and such familiarity with the best astronomical writings in a man whose professional duties are of an entirely different kind.
